\section*{Chapitre \ref{ch:b1:stereochemistry-in-depth} --- La stéréochimie approfondie}

\begin{solution}{exo:b1:stereochemistry-in-depth:1}
\omterm{def:b1:stereochemistry-in-depth:isomers}{Isomères de constitution}~; \omterm{def:b1:stereochemistry-in-depth:enantiomers}{énantiomères}~; \omterm{def:b1:stereochemistry-in-depth:enantiomers}{diastéréoisomères}
(\omterm{def:b1:stereochemistry-in-depth:isomers}{stéréoisomères} qui ne sont pas images l'un de l'autre)~; deux
\omterm{def:b1:stereochemistry-in-depth:isomers}{conformations} d'un même composé~; \omterm{def:b1:stereochemistry-in-depth:enantiomers}{diastéréoisomères} (\cip{R,S} est la
forme méso).
\end{solution}

\begin{solution}{exo:b1:stereochemistry-in-depth:2}
\ce{-OH} $>$ \ce{-NH2} $>$ \ce{-CH3} $>$ \ce{-H} (O, N, C, H).
\ce{-COOH} (O,O,O) $>$ \ce{-CHO} (O,O,H) $>$ \ce{-CH2OH} (O,H,H) $>$
\ce{-CH(CH3)2} (C,C,H). \ce{-Br} $>$ \ce{-Cl} $>$ \ce{-CCl3} (Cl,Cl,Cl)
$>$ \ce{-CH2Cl} (Cl,H,H)~: le premier atome décide avant ses voisins.
\end{solution}

\begin{solution}{exo:b1:stereochemistry-in-depth:3}
2,3-Dichlorobutane~: deux centres équivalents, \cip{R,R}, \cip{S,S} et
la forme méso \cip{R,S}~: trois. 2-Bromo-3-chlorobutane~: deux centres
différents, quatre \omterm{def:b1:stereochemistry-in-depth:isomers}{stéréoisomères}, pas de forme méso.
Pentane-2,4-diol~: trois, \cip{2R,4S} étant méso.
\end{solution}

\begin{solution}{exo:b1:stereochemistry-in-depth:4}
Sur C5~: isopropényle (C,C,C, la liaison double comptant deux fois) $>$
C6 $>$ C4
$>$ H. C6 et C4 portent tous deux (C,H,H)~; un rang plus loin, C6 mène
au carbone du carbonyle (O,O,C) et C4 à C3 (C,C,H)~: C6 l'emporte. Le
groupe isopropényle étant en avant et H en arrière, l'ordre
isopropényle $\to$ C6 $\to$ C4 tourne dans le sens horaire tel que
dessiné~: \cip{R}, la carvone de la menthe verte~; celle au coin hachuré
est \cip{S}.
\end{solution}

\begin{solution}{exo:b1:stereochemistry-in-depth:5}
$\theta = 0$~: méthyles éclipsés, \qty{16.5}{kJ/mol}~; $60^\circ$~:
gauche, environ
\qty{2.9}{kJ/mol} (minimum 2.8 à $62^\circ$)~; $120^\circ$~: méthyle
éclipsant H,
\qty{15.2}{kJ/mol}~; $180^\circ$~: anti, 0. À énergies égales, deux
\omterm{def:b1:stereochemistry-in-depth:conformations}{conformations gauches} pour une anti~: deux tiers de gauches.
\end{solution}

\begin{solution}{exo:b1:stereochemistry-in-depth:6}
Chaîne verticale avec \ce{COOH} en haut et \ce{CH3} en bas~; rangs
$\ce{NH2} > \ce{COOH} > \ce{CH3} > \ce{H}$. En plaçant \ce{NH2} à gauche
et H à droite, $\ce{NH2} \to \ce{COOH} \to \ce{CH3}$ tourne dans le sens
horaire tel qu'on le voit~; H est horizontal (vers le lecteur), donc le
descripteur est inversé~: \cip{S}. Le groupe amino est à gauche.
\end{solution}

\begin{solution}{exo:b1:stereochemistry-in-depth:7}
\textit{cis}~: sur des carbones adjacents, une liaison vers le haut et
une vers le bas parmi les positions axiales et \omterm{def:b1:stereochemistry-in-depth:conformations}{équatoriales} donnent
axial--équatorial dans les deux \omterm{def:b1:stereochemistry-in-depth:conformations}{chaises}.
\textit{trans}~: diéquatorial ou diaxial. Dans la \omterm{def:b1:stereochemistry-in-depth:conformations}{chaise} diéquatoriale,
les deux méthyles sont gauches l'un par rapport à l'autre (une
interaction)~; l'isomère \textit{cis} a cette interaction plus un méthyle
\omterm{def:b1:stereochemistry-in-depth:conformations}{axial}, environ \qty{5.7}{kJ/mol} de plus. L'isomère \textit{trans} est
le plus stable.
\end{solution}

\begin{solution}{exo:b1:stereochemistry-in-depth:8}
$[\alpha] = 13.2/(2.00 \times 0.100) = +66.0^\circ$. Tube de
\qty{1.00}{dm}~: $+6.60^\circ$. L'énantiomère~: $-13.2^\circ$ dans le
tube de \qty{2.00}{dm}.
\end{solution}

\begin{solution}{exo:b1:stereochemistry-in-depth:9}
\qty{25}{\celsius}~: $\exp(5700/(8.314 \times 298.15)) = 10.0$, fraction
$10.0/11.0 = \qty{91}{\percent}$. \qty{-50}{\celsius}~: $\exp(5700/(8.314
\times 223.15)) = 21.6$, \qty{96}{\percent}.
\end{solution}

\begin{solution}{exo:b1:stereochemistry-in-depth:10}
Dans l'éthane, les six hydrogènes sont semblables~: toutes les
\omterm{def:b1:stereochemistry-in-depth:conformations}{conformations décalées} sont identiques. Dans le butane, le méthyle
arrière peut se placer à l'opposé du méthyle avant (anti) ou à côté de
lui ($\pm 60^\circ$, gauche), où les deux méthyles s'encombrent.
Populations~: anti 1, gauche $2\exp(-2.8/2.48) = 0.65$~: environ
\qty{61}{\percent} d'anti et \qty{39}{\percent} de gauche.
\end{solution}

\begin{solution}{exo:b1:stereochemistry-in-depth:11}
C2 et C4 sont stéréogènes. C3 porte H, OH et deux branches identiques
par leur constitution~; elles ne diffèrent que lorsque C2 et C4 ont des
descripteurs opposés, et c'est seulement alors que C3 est stéréogène
(pseudo-asymétrique). Isomères~: \cip{2R,4R} et \cip{2S,4S}, un couple
d'\omterm{def:b1:stereochemistry-in-depth:enantiomers}{énantiomères} (C3 non stéréogène)~; \cip{2R,4S} avec les deux
dispositions en C3, deux formes méso. Quatre en tout.
\end{solution}

\begin{solution}{exo:b1:stereochemistry-in-depth:12}
$[\alpha]_{\text{éch}} = 2.31/(1.00 \times 0.050) = +46.2^\circ$, et
$2x - 1 = 46.2/66.0 = 0.700$~: $x = 0.850$, \qty{85}{\percent} de
l'énantiomère $(+)$ et \qty{15}{\percent} du $(-)$.
\end{solution}

\begin{solution}{pb:b1:stereochemistry-in-depth:1}
\textbf{1.} Cycle du cyclohexane~: C1 porte OH, C2 le groupe isopropyle,
C5 le méthyle~; C1, C2 et C5 sont stéréogènes.
\textbf{2.} $2^3 = 8$, quatre couples d'\omterm{def:b1:stereochemistry-in-depth:enantiomers}{énantiomères}.
\textbf{3.} Aucune disposition ne présente de plan de symétrie~: les
trois substituants sont tous différents.
\textbf{4.} \ce{OH} $>$ C2 (C,C,H) $>$ C6 (C,H,H) $>$ H.
\textbf{5.} C1 (O,C,H) $>$ isopropyle (C,C,H) $>$ C3 (C,H,H) $>$ H.
\textbf{6.} C6 et C4 sont tous deux (C,H,H)~; ensuite, C6 mène à C1
(O,C,H), C4 à C3
(C,H,H)~: C6 $>$ C4 $>$ \ce{CH3} $>$ H.
\textbf{7.} \cip{1S,2R,5S}.
\textbf{8.} Dans une \omterm{def:b1:stereochemistry-in-depth:conformations}{chaise}, deux liaisons \omterm{def:b1:stereochemistry-in-depth:conformations}{équatoriales} sur des carbones
adjacents (C1, C2) pointent l'une vers le haut, l'autre vers le bas~:
\textit{trans}~; sur les carbones 1 et 3 du cycle (C1 et C5, en passant
par C6), les deux liaisons \omterm{def:b1:stereochemistry-in-depth:conformations}{équatoriales} pointent du même côté~:
\textit{cis}. C'est la disposition du menthol.
\textbf{9.} Les trois deviennent \omterm{def:b1:stereochemistry-in-depth:conformations}{axiaux}.
\textbf{10.} Dans la \omterm{def:b1:stereochemistry-in-depth:conformations}{chaise} inversée, le méthyle seul coûterait environ
\qty{5.7}{kJ/mol} et le groupe isopropyle, plus gros, davantage~; le
groupe OH ajoute un coût propre plus faible, et le OH et le méthyle
\omterm{def:b1:stereochemistry-in-depth:conformations}{axiaux}, portés par les carbones 1 et 3 et situés sur la même face,
s'encombreraient aussi. On dépasse largement \qty{12}{kJ/mol}, soit un
rapport supérieur à $\exp(12000/2479) \approx 130$ contre un~: le
menthol est presque entièrement dans la \omterm{def:b1:stereochemistry-in-depth:conformations}{chaise} tout \omterm{def:b1:stereochemistry-in-depth:conformations}{équatoriale}.
\textbf{11.} Un diastéréoisomère~: un seul des trois centres diffère.
\textbf{12.} L'isopropyle et le méthyle étant \omterm{def:b1:stereochemistry-in-depth:conformations}{équatoriaux}, le OH est
\omterm{def:b1:stereochemistry-in-depth:conformations}{axial}.
\textbf{13.} Les \omterm{def:b1:stereochemistry-in-depth:enantiomers}{diastéréoisomères} diffèrent par toutes leurs propriétés
(ici les \omterm{def:b1:intermolecular-forces-solvents:hbond}{liaisons hydrogène} d'un OH \omterm{def:b1:stereochemistry-in-depth:conformations}{axial} ou \omterm{def:b1:stereochemistry-in-depth:conformations}{équatorial})~; les
\omterm{def:b1:stereochemistry-in-depth:enantiomers}{énantiomères} ne diffèrent que vis-à-vis de partenaires chiraux.
\textbf{14.} $[\alpha] = -2.46/(1.00 \times 0.0500) = -49.2^\circ$, dans
l'intervalle du manuel.
% ledger: alpha:l-menthol-range
\textbf{15.} $+49.2^\circ$~; $0$.
\textbf{16.} $-1.97/0.0500 = -39.4^\circ$.
\textbf{17.} $\alpha = \ell c[\alpha]_{(-)}\bigl(x - (1 - x)\bigr) =
\ell c[\alpha]_{(-)}(2x - 1)$.
\textbf{18.} $x = \frac12(1 + \alpha/\alpha_{\text{réf}})$ à $\ell$ et $c$
égaux.
\textbf{19.} Oui~: tout autre composé optiquement actif ajoute son propre
terme (\omterm{prop:b1:stereochemistry-in-depth:biot}{loi de Biot}). Il faut vérifier la pureté chimique par une autre
méthode (chromatographie) avant de lire le pouvoir rotatoire comme une
composition.
\textbf{20.} $x = \frac12(1 + 1.97/2.46) = 0.900$.
\textbf{21.} $x = \frac12(1 + 2.40/2.46) = 0.988$.
\textbf{22.} Le second (\qty{98.8}{\percent}).
\textbf{23.} $-3.94^\circ$~; la même erreur de lecture représente alors
une fraction plus petite de l'angle.
\textbf{24.} \textbf{\qty{90.0}{\percent} de $(-)$-menthol} (et
\qty{10.0}{\percent} de son énantiomère).
\end{solution}
